\documentclass[10pt,letterpaper]{article}
\usepackage[utf8]{inputenc}
\usepackage[intlimits]{amsmath}
\usepackage{amsfonts}
\usepackage{amssymb}
\usepackage{ragged2e}
\usepackage[letterpaper, margin=0.5in]{geometry}
\usepackage{graphicx}
\usepackage{cancel}
\usepackage{mathtools}
\usepackage{tabularx}
\usepackage{arydshln}
\usepackage{tensor}
\usepackage{array}
\usepackage{xcolor}
\usepackage[boxed]{algorithm}
\usepackage[noend]{algpseudocode}
\usepackage{listings}
\usepackage{textcomp}
% \usepackage[pdf,tmpdir,singlefile]{graphviz}
\usepackage{mathrsfs}
\usepackage{bbm}
\usepackage{tikz}
\usepackage{tikz-cd}
\usepackage{enumitem}
\usepackage{arydshln}
\usepackage{relsize}
\usepackage{multirow,multicol}
\usepackage{scalerel}
\usepackage{upgreek}
\usepackage{ifthen}

\usetikzlibrary{bayesnet}
\setlist{noitemsep}

%%%%%%%%%%%%%%%%%%%%%%%%%%%%%
% Formatting commands
%%%%%%%%%%%%%%%%%%%%%%%%%%%%%
\newcommand{\n}{\hfill\break}
\newcommand{\nn}{\vspace{0.5\baselineskip}\n}
\newcommand{\up}{\vspace{-\baselineskip}}
\newcommand{\hangblock}[2]{\par\noindent\settowidth{\hangindent}{\textbf{#1: }}\textbf{#1: }\nolinebreak#2}
\newcommand{\lemma}[1]{\hangblock{Lemma}{#1}}
\newcommand{\defn}[1]{\hangblock{Defn}{#1}}
\newcommand{\thm}[1]{\hangblock{Thm}{#1}}
\newcommand{\cor}[1]{\hangblock{Cor}{#1}}
\newcommand{\prop}[1]{\hangblock{Prop}{#1}}
\newcommand{\ex}[1]{\hangblock{Ex}{#1}}
\newcommand{\exer}[1]{\hangblock{Exer}{#1}}
\newcommand{\fact}[1]{\hangblock{Fact}{#1}}
\newcommand{\remark}[1]{\hangblock{Remark}{#1}}
\newcommand{\proven}{\;$\square$\n}
\newcommand{\problem}[1]{\par\noindent{\nolinebreak#1}\n}
\newcommand{\problempart}[2]{\par\noindent\indent{}\settowidth{\hangindent}{\textbf{(#1)} \indent{}}\textbf{(#1) }\nolinebreak#2\n}
\newcommand{\ptxt}[1]{\textrm{\textnormal{#1}}}
\newcommand{\inlineeq}[1]{\centerline{$\displaystyle #1$}}
\newcommand{\pageline}{\up\par\noindent\rule{\textwidth}{0.1pt}}

%%%%%%%%%%%%%%%%%%%%%%%%%%%%%
% Math commands
%%%%%%%%%%%%%%%%%%%%%%%%%%%%%
% Set Theory
\newcommand{\card}[1]{\left|#1\right|}
\newcommand{\set}[1]{\left\{#1\right\}}
\newcommand{\setmid}{\;\middle|\;}
\newcommand{\ps}[1]{\mathcal{P}\left(#1\right)}
\newcommand{\pfinite}[1]{\mathcal{P}^{\ptxt{finite}}\left(#1\right)}
\newcommand{\naturals}{\mathbb{N}}
\newcommand{\N}{\naturals}
\newcommand{\integers}{\mathbb{Z}}
\newcommand{\Z}{\integers}
\newcommand{\rationals}{\mathbb{Q}}
\newcommand{\Q}{\rationals}
\newcommand{\reals}{\mathbb{R}}
\newcommand{\R}{\reals}
\newcommand{\complex}{\mathbb{C}}
\newcommand{\C}{\complex}
\newcommand{\halfPlane}{\mathbb{H}}
\let\HSym\H
\let\H\relax
\newcommand{\H}{\halfPlane}
\newcommand{\comp}{^{\complement}}
\DeclareMathOperator{\Hom}{Hom}
\newcommand{\Ind}{\mathbbm{1}}
\newcommand{\cut}{\setminus}
\DeclareMathOperator{\elem}{elem}

% Differentiable Manifolds
\newcommand{\RP}{\mathbb{RP}}
\newcommand{\CP}{\mathbb{RP}}
\newcommand{\osubset}{\overset{\mathclap{\scalebox{0.5}{\ptxt{open}}}}{\subset}}
\newcommand{\osubseteq}{\overset{\mathclap{\scalebox{0.5}{\ptxt{open}}}}{\subseteq}}
\newcommand{\osupset}{\overset{\mathclap{\scalebox{0.5}{\ptxt{open}}}}{\supset}}
\newcommand{\osupseteq}{\overset{\mathclap{\scalebox{0.5}{\ptxt{open}}}}{\supseteq}}
\newcommand{\pdat}[3]{\left.\pd{#1}{#2}\right|_{#3}}
\DeclareMathOperator{\so}{so}
\DeclareMathOperator{\codim}{codim}
\DeclareMathOperator{\Diff}{Diff}
\let\dSym\d
\let\d\relax
\newcommand{\d}{\partial}
\DeclareMathOperator{\gl}{gl}
\DeclareMathOperator{\Ad}{Ad}
\newcommand{\comm}[1]{\left[#1\right]}

% Graph Theory
\let\deg\relax
\DeclareMathOperator{\deg}{deg}
\newcommand{\degp}{\ptxt{deg}^{+}}
\newcommand{\degn}{\ptxt{deg}^{-}}
\newcommand{\precdot}{\mathrel{\ooalign{$\prec$\cr\hidewidth\hbox{$\cdot\mkern0.5mu$}\cr}}}
\newcommand{\succdot}{\mathrel{\ooalign{$\cdot\mkern0.5mu$\cr\hidewidth\hbox{$\succ$}\cr\phantom{$\succ$}}}}
\DeclareMathOperator{\cl}{cl}
\DeclareMathOperator{\affdim}{affdim}

% Probability
\newcommand{\parSymbol}{\P}
\newcommand{\Prob}{\mathbb{P}}
\renewcommand{\P}{\Prob}
\newcommand{\Avg}{\mathbb{E}}
\newcommand{\E}{\Avg}
\DeclareMathOperator{\Var}{Var}
\DeclareMathOperator{\cov}{cov}
\DeclareMathOperator{\Unif}{Unif}
\DeclareMathOperator{\Binom}{Binom}
\newcommand{\CI}{\mathrel{\text{\scalebox{1.07}{$\perp\mkern-10mu\perp$}}}}
\DeclareMathOperator{\Ber}{Ber}
\DeclareMathOperator{\Bin}{Bin}
\DeclareMathOperator{\Geom}{Geom}
\DeclareMathOperator{\Poisson}{Poisson}
\DeclareMathOperator{\Exp}{Exp}
\DeclareMathOperator{\Beta}{Beta}
\DeclareMathOperator{\Normal}{N}

% Standard Math
\newcommand{\inv}{^{-1}}
\newcommand{\abs}[1]{\left|#1\right|}
\newcommand{\ceil}[1]{\left\lceil{}#1\right\rceil{}}
\newcommand{\floor}[1]{\left\lfloor{}#1\right\rfloor{}}
\newcommand{\conj}[1]{\overline{#1}}
\newcommand{\of}{\circ}
\newcommand{\tri}{\triangle}
\newcommand{\inj}{\hookrightarrow}
\newcommand{\surj}{\twoheadrightarrow}
\newcommand{\ndiv}{\nmid}
\renewcommand{\epsilon}{\varepsilon}
\newcommand{\divides}{\mid}
\newcommand{\ndivides}{\nmid}
\DeclareMathOperator{\lcm}{lcm}
\DeclareMathOperator{\sgn}{sgn}
\newcommand{\map}[4]{\!\!\!\begin{array}[t]{rcl}#1 & \!\!\!\!\to & \!\!\!\!#2\\ {}#3 & \!\!\!\!\mapsto & \!\!\!\!#4\end{array}}
\newcommand{\bigsum}[2]{\smashoperator[lr]{\sum_{\scalebox{#1}{$#2$}}}}
\DeclareMathOperator{\gcf}{gcf}
\newcommand{\restr}[2]{\left.#1\right|_{#2}}

% Linear Algebra
\newcommand{\Id}{\textrm{\textnormal{Id}}}
\newcommand{\im}{\textrm{\textnormal{im}}}
\newcommand{\norm}[1]{\abs{\abs{#1}}}
\newcommand{\tpose}{^{T}\!}
\newcommand{\iprod}[1]{\left<#1\right>}
\newcommand{\giprod}{\iprod{\;\,,\;}}
\DeclareMathOperator{\tr}{tr}
\DeclareMathOperator{\trace}{tr}
\newcommand{\chgBasMat}[3]{\!\!\tensor*[_{#1}]{\left[#2\right]}{_{#3}}}
\newcommand{\vecBas}[2]{\tensor*[]{\left[#1\right]}{_{#2}}}
\DeclareMathOperator{\GL}{GL}
\DeclareMathOperator{\Mat}{Mat}
\DeclareMathOperator{\vspan}{span}
\DeclareMathOperator{\rank}{rank}
\newcommand{\V}[1]{\vec{#1}}
\DeclareMathOperator{\proj}{proj}
\DeclareMathOperator{\compProj}{comp}
\DeclareMathOperator{\row}{row}
\newcommand{\smallPMatrix}[1]{\paren{\begin{smallmatrix}#1\end{smallmatrix}}}
\newcommand{\smallBMatrix}[1]{\brack{\begin{smallmatrix}#1\end{smallmatrix}}}
\newcommand{\pmat}[1]{\begin{pmatrix}#1\end{pmatrix}}
\newcommand{\bmat}[1]{\begin{bmatrix}#1\end{bmatrix}}
\newcommand{\dual}{^{*}}
\newcommand{\pinv}{^{\dagger}}
\newcommand{\horizontalMatrixLine}{\ptxt{\rotatebox[origin=c]{-90}{$|$}}}
\DeclareMathOperator{\range}{range}
\DeclareMathOperator{\Symm}{Symm}
\DeclareMathOperator{\SU}{SU}
\DeclareMathOperator{\U}{U}

% Multilinear Algebra
\let\Lsym\L
\let\L\relax
\DeclareMathOperator{\L}{\mathscr{L}}
\DeclareMathOperator{\A}{\mathcal{A}}
\DeclareMathOperator{\Alt}{Alt}
\DeclareMathOperator{\Sym}{Sym}
\newcommand{\ot}{\otimes}
\newcommand{\ox}{\otimes}
\DeclareMathOperator{\asc}{asc}
\DeclareMathOperator{\asSet}{set}
\DeclareMathOperator{\sort}{sort}
\DeclareMathOperator{\ringA}{\mathring{A}}
\DeclareMathOperator{\Sh}{Sh}
\DeclareMathOperator{\Bil}{Bil}

% Topology
\newcommand{\closure}[1]{\overline{#1}}
\newcommand{\uball}{\mathcal{U}}
\DeclareMathOperator{\Int}{Int}
\DeclareMathOperator{\Ext}{Ext}
\DeclareMathOperator{\Bd}{Bd}
\DeclareMathOperator{\rInt}{rInt}
\DeclareMathOperator{\ch}{ch}
\DeclareMathOperator{\ah}{ah}
\newcommand{\LargerTau}{\mathlarger{\mathlarger{\mathlarger{\mathlarger{\tau}}}}}
\newcommand{\Tau}{\mathcal{T}}

% Analysis
\DeclareMathOperator{\Graph}{Graph}
\DeclareMathOperator{\epi}{epi}
\DeclareMathOperator{\hypo}{hypo}
\DeclareMathOperator{\supp}{supp}
\newcommand{\lint}[2]{\underset{#1}{\overset{#2}{{\color{black}\underline{{\color{white}\overline{{\color{black}\int}}\color{black}}}}}}}
\newcommand{\uint}[2]{\underset{#1}{\overset{#2}{{\color{white}\underline{{\color{black}\overline{{\color{black}\int}}\color{black}}}}}}}
\newcommand{\alignint}[2]{\underset{#1}{\overset{#2}{{\color{white}\underline{{\color{white}\overline{{\color{black}\int}}\color{black}}}}}}}
\newcommand{\extint}{\ptxt{ext}\int}
\newcommand{\extalignint}[2]{\ptxt{ext}\alignint{#1}{#2}}
\newcommand{\conv}{\ast}
\newcommand{\pd}[2]{\frac{\partial{}#1}{\partial{}#2}}
\newcommand{\del}{\nabla}
\DeclareMathOperator{\grad}{grad}
\DeclareMathOperator{\curl}{curl}
\let\div\relax
\DeclareMathOperator{\div}{div}
\DeclareMathOperator{\vol}{vol}

% Complex Analysis
\let\Re\relax
\DeclareMathOperator{\Re}{Re}
\let\Im\relax
\DeclareMathOperator{\Im}{Im}
\DeclareMathOperator{\Res}{Res}

% Abstract Algebra
\DeclareMathOperator{\ord}{ord}
\newcommand{\generated}[1]{\left<#1\right>}
\newcommand{\cycle}[1]{\smallPMatrix{#1}}
\newcommand{\id}{\ptxt{id}}
\newcommand{\iso}{\cong}
\DeclareMathOperator{\Aut}{Aut}
\DeclareMathOperator{\SL}{SL}
\DeclareMathOperator{\op}{op}
\newcommand{\isom}[4]{\!\!\!\begin{array}[t]{rcl}#1 & \!\!\!\!\overset{\sim}{\to} & \!\!\!\!#2\\ #3 & \!\!\!\!\mapsto & \!\!\!\!#4\end{array}}
\newcommand{\F}{\mathbb{F}}
\newcommand{\acts}{\;\reflectbox{\rotatebox[origin=c]{-90}{$\circlearrowright$}}\;}
\newcommand{\disjunion}{\mathrel{\text{\scalebox{1.07}{$\perp\mkern-10mu\perp$}}}}
\DeclareMathOperator{\SO}{SO}
\DeclareMathOperator{\stab}{stab}
\DeclareMathOperator{\nullity}{nullity}
\DeclareMathOperator{\Perm}{Perm}
\DeclareMathOperator{\nsubgp}{\vartriangleleft}
\DeclareMathOperator{\notnsubgp}{\ntriangleleft}
\newcommand{\presentation}[2]{\left<#1\;\middle|\;#2\right>}
\DeclareMathOperator{\Char}{char}
\DeclareMathOperator{\fchar}{char}
\DeclareMathOperator{\triv}{triv}
\DeclareMathOperator{\reg}{reg}
\DeclareMathOperator{\std}{std}
\DeclareMathOperator{\Func}{Func}
\DeclareMathOperator{\End}{End}

% Convex Optimization
\let\sectionSymbol\S
\let\S\relax
\newcommand{\S}{\mathbb{S}}
\DeclareMathOperator{\dist}{dist}
\DeclareMathOperator{\dom}{dom}
\DeclareMathOperator{\diag}{diag}
\DeclareMathOperator{\ones}{\mathbbm{1}}
\newcommand{\minimizeOver}[3]{\begin{array}{rl}\underset{#1}{\ptxt{minimize}} & #2\\ \ptxt{subject to} & #3\end{array}}
\newcommand{\maximizeOver}[3]{\begin{array}{rl}\underset{#1}{\ptxt{maximize}} & #2\\ \ptxt{subject to} & #3\end{array}}
\newcommand{\minimizationProblem}[2]{\minimizeOver{}{#1}{#2}}
\newcommand{\maximizationProblem}[2]{\maximizeOver{}{#1}{#2}}
\newcommand{\minimizeOverUnconstrained}[2]{\begin{array}{rl}\underset{#1}{\ptxt{minimize}} & #2\end{array}}
\newcommand{\maximizeOverUnconstrained}[2]{\begin{array}{rl}\underset{#1}{\ptxt{maximize}} & #2\end{array}}
\newcommand{\minimizationUnconstrained}[1]{\minimizeOverUnconstrained{}{#1}}
\newcommand{\maximizationUnconstrained}[1]{\maximizeOverUnconstrained{}{#1}}
\DeclareMathOperator{\argmin}{argmin}
\DeclareMathOperator{\argmax}{argmax}

% Proofs
\newcommand{\st}{s.t.}
\newcommand{\unique}{!}
\newcommand{\iffdef}{\overset{\ptxt{def}}{\Leftrightarrow}}
\newcommand{\eqVertical}{\rotatebox[origin=c]{90}{=}}
\newcommand{\mapsfrom}{\mathrel{\reflectbox{\ensuremath{\mapsto}}}}
\newcommand{\mapsdown}{\rotatebox[origin=c]{-90}{$\mapsto$}\mkern2mu}
\newcommand{\mapsup}{\rotatebox[origin=c]{90}{$\mapsto$}\mkern2mu}
\newcommand{\from}{\!\mathrel{\reflectbox{\ensuremath{\to}}}}
\newcommand{\labeledeq}[1]{\overset{\mathclap{\ptxt{#1}}}{=}}
\newcommand{\eqdef}{\labeledeq{def}}

% Brackets
\newcommand{\paren}[1]{\left(#1\right)}
\renewcommand{\brack}[1]{\left[#1\right]}
\renewcommand{\brace}[1]{\left\{#1\right\}}
\newcommand{\ang}[1]{\left<#1\right>}

% Algorithms
\algrenewcommand{\algorithmiccomment}[1]{\hskip 1em \texttt{// #1}}
\algrenewcommand\algorithmicrequire{\textbf{Input:}}
\algrenewcommand\algorithmicensure{\textbf{Output:}}
\newcommand{\algP}{\ptxt{\textbf{P}}}
\newcommand{\algNP}{\ptxt{\textbf{NP}}}
\newcommand{\algNPC}{\ptxt{\textbf{NP-Complete}}}
\newcommand{\algNPH}{\ptxt{\textbf{NP-Hard}}}
\newcommand{\algEXP}{\ptxt{\textbf{EXP}}}
\DeclareMathOperator{\fl}{fl}

%%%%%%%%%%%%%%%%%%%%%%%%%%%%%
% Other commands and shorthand
%%%%%%%%%%%%%%%%%%%%%%%%%%%%%
\newcommand{\flag}[1]{\textbf{\textcolor{red}{#1}}}
\let\uSym\u
\let\u\relax
\newcommand{\u}[1]{\underline{#1}}
\let\bSym\b
\let\b\relax
\newcommand{\b}[1]{\textbf{#1}}
\let\iSym\i
\let\i\relax
\newcommand{\i}[1]{\textit{#1}}
\let\scSym\sc
\let\sc\relax
\newcommand{\sc}[1]{\textsc{#1}}
\newcommand{\mf}[1]{\mathfrak{#1}}
\newcommand{\mc}[1]{\mathcal{#1}}

%%%%%%%%%%%%%%%%%%%%%%%%%%%%%%%%%%%%%%%
% Make l's curvy in math environments %
%%%%%%%%%%%%%%%%%%%%%%%%%%%%%%%%%%%%%%%
\mathcode`l="8000
\begingroup
\makeatletter
\lccode`\~=`\l
\DeclareMathSymbol{\lsb@l}{\mathalpha}{letters}{`l}
\lowercase{\gdef~{\ifnum\the\mathgroup=\m@ne \ell \else \lsb@l \fi}}%
\endgroup

%%%%%%%%%%%%%%%%%%%%%%%%%
% Fix \vdots and \ddots %
%%%%%%%%%%%%%%%%%%%%%%%%%
\usepackage{letltxmacro}
\LetLtxMacro\orgvdots\vdots
\LetLtxMacro\orgddots\ddots

\makeatletter
\DeclareRobustCommand\vdots{%
	\mathpalette\@vdots{}%
}
\newcommand*{\@vdots}[2]{%
	% #1: math style
	% #2: unused
	\sbox0{$#1\cdotp\cdotp\cdotp\m@th$}%
	\sbox2{$#1.\m@th$}%
	\vbox{%
		\dimen@=\wd0 %
		\advance\dimen@ -3\ht2 %
		\kern.5\dimen@
		% remove side bearings
		\dimen@=\wd2 %
		\advance\dimen@ -\ht2 %
		\dimen2=\wd0 %
		\advance\dimen2 -\dimen@
		\vbox to \dimen2{%
			\offinterlineskip
			\copy2 \vfill\copy2 \vfill\copy2 %
		}%
	}%
}
\DeclareRobustCommand\ddots{%
	\mathinner{%
		\mathpalette\@ddots{}%
		\mkern\thinmuskip
	}%
}
\newcommand*{\@ddots}[2]{%
	% #1: math style
	% #2: unused
	\sbox0{$#1\cdotp\cdotp\cdotp\m@th$}%
	\sbox2{$#1.\m@th$}%
	\vbox{%
		\dimen@=\wd0 %
		\advance\dimen@ -3\ht2 %
		\kern.5\dimen@
		% remove side bearings
		\dimen@=\wd2 %
		\advance\dimen@ -\ht2 %
		\dimen2=\wd0 %
		\advance\dimen2 -\dimen@
		\vbox to \dimen2{%
			\offinterlineskip
			\hbox{$#1\mathpunct{.}\m@th$}%
			\vfill
			\hbox{$#1\mathpunct{\kern\wd2}\mathpunct{.}\m@th$}%
			\vfill
			\hbox{$#1\mathpunct{\kern\wd2}\mathpunct{\kern\wd2}\mathpunct{.}\m@th$}%
		}%
	}%
}
\makeatother

\newcommand{\B}{
	\begin{tikzpicture}
	\filldraw [fill=red, draw=black] (0, 0) rectangle (0.37, 0.45);
	\draw [line width=0.5mm, white ] (0.1,0.08) -- (0.1,0.38);
	\draw[line width=0.5mm, white ] (0.1, 0.35) .. controls (0.2, 0.35) and (0.4, 0.2625) .. (0.1, 0.225);
	\draw[line width=0.5mm, white ] (0.1, 0.225) .. controls (0.2, 0.225) and (0.4, 0.1625) .. (0.1, 0.1);
	\end{tikzpicture}
}

% Allow custom symbols as arrows in tikz-cd
\tikzset{
  symbol/.style={
    draw=none,
    every to/.append style={
      edge node={node [sloped, allow upside down, auto=false]{$#1$}}}
  }
}
% Example: \arrow[r,symbol=\cong]
% https://tex.stackexchange.com/questions/394154/how-to-include-inclusion-subgroup-relationship-in-tikz-cd-diagram

\author{Thomas Cohn}
\title{Stats 426 Lecture 4}
\date{2/1/21} % Can also use \today

\begin{document}
\maketitle
\setlength\RaggedRightParindent{\parindent}
\RaggedRight

\subsection*{Continuous Random Variables and Probability Distributions}

\par\noindent
Recall: A random variable $X$ is continuous if its set of possible values contains an entire interval of numbers. The distribution of $X$ is characterized by a pdf $f(x)$ \st{} $\forall{}a,b\in\R$ with $a<b$, $P(a<X<b)=\int_{a}^{b}f(x)\,dx$. Also, $f(x)\ge{}0$ $\forall{}x\in\R$, $f$ is piecewise continuous, and $\int_{\R}f=1$.\n

\defn{
	Suppose $X$ is a continuous random variable with density function $f$. Then the \u{cumulative distribution function} (or \u{cdf}) of $X$ is
	\[
		F(x)\eqdef{}P(X\le{}x)=\int_{-\infty}^{x}f(u)\,du
	\]
}

\par\noindent
Observe: $\lim_{s\to-\infty}F(s)=0$, $\lim_{s\to\infty}F(s)=1$.\n

\par\noindent
Some facts about $X$, an arbitrary continuous random variable:
\begin{itemize}
	\item $P(X=c)=0$, $\forall{}c\in\R$
	\item $P(a\le{}X\le{}b)=F(b)-F(a)$
	\item $P(a<X<b)=P(a<X\le{}b)=P(a\le{}X<b)=P(a\le{}X\le{}b)$
	\item $\forall{}x$ where $F'(x)$ exists, $\frac{d}{dx}F(x)=f(x)$.
\end{itemize}

\subsection*{Quantiles}

\defn{
	If a cdf $F$ is strictly increasing, for $p\in(0,1)$, $x=F\inv(p)$ is called the \u{$p$th quantile}.\n
}

\defn{
	The $0.5$ quantile is the \u{median}.\n
	The $0.25$ quantile is the \u{lower quantile}.\n
	The $0.75$ quantile is the \u{upper quantile}.\n
}

\defn{
	$X$ follows an \u{exponential distribution} with parameter $\lambda$ ($X\sim\Exp(\lambda)$) if
	\[
		f(x;\lambda)=\left\{\begin{array}{ll}
			\lambda{}e^{-\lambda{}x} & x\ge{}0\\
			0 & x<0
		\end{array}\right.\qquad
		F(x;\lambda)=\left\{\begin{array}{ll}
			1-e^{-\lambda{}x} & x\ge{}0\\
			0 & x<0\\
		\end{array}\right.
	\]
	Note: We must have $\lambda>0$.\n
}

\par\noindent
$\Exp$ is often used as the distribution of lifetimes, or times between the occurrence of successive events.\n

\par\noindent
Suppose $X\sim\Exp(\lambda)$, and $x,x_{0}>0$. Then $P(X\ge{}x+x_{0}\mid{}X\ge{}x_{0})=P(X\ge{}x)$.\nn
Proof:
\[
	P(X\ge{}x+x_{0}\mid{}X\ge{}x_{0})=\frac{P(X\ge{}x+x_{0},X\ge{}x_{0})}{P(X\ge{}x_{0})}=\frac{P(X\ge{}x+x_{0})}{P(X\ge{}x_{0})}=\frac{1-(1-e^{-\lambda{}(x+x_{0})})}{1-(1-e^{-\lambda{}x_{0}})}=e^{-\lambda{}x}=P(X\ge{}x)
\]
\proven

\ex{
	The lifetime of a component is exponentially-distributed, with parameter $\lambda=3$. If it has already worked for $10$ hours, what is the probability that it works for $4$ more hours?
	\[
		P(X\ge{}10+4\mid{}X\ge{}10)=P(X\ge{}4)=e^{-3\cdot{}4}=e^{-12}
	\]
}

\defn{
	$X$ follows a \u{gamma distribution} with \u{shape parameter} $\alpha>0$ and \u{scale parameter} $\lambda>0$ ($X\sim\gamma(\lambda,\alpha)$) if
	\[
		f(x)=\left\{\begin{array}{ll}
			\frac{\lambda^{\alpha}}{\Gamma(\alpha)}x^{\alpha-1}e^{-\lambda{}x} & x\ge{}0\\
			0 & x<0
		\end{array}\right.
	\]
	where the \u{gamma function} $\Gamma$ satisfies
	\[
		\Gamma(\alpha)=\int_{0}^{\infty}x^{\alpha-1}e^{-x}\,dx
	\]
}

\par\noindent
Some properties of the gamma function:
\begin{itemize}
	\item $\forall\alpha>1$, $\Gamma(\alpha)=(\alpha-1)\Gamma(\alpha-1)$
	\item $\forall{}n\in\N$, $\Gamma(n)=(n-1)!$
	\item $\Gamma(\frac{1}{2})=\sqrt{\pi}$
\end{itemize}

\par\noindent
Some properties of the gamma distribution:
\begin{itemize}
	\item $\Exp$ is a special case with $\alpha=1$.
	\item The sum of i.i.d. exponential random variables follows a gamma distribution. If $X_{1},\ldots,X_{n}\overset{\ptxt{i.i.d.}}{\sim}\Exp(\lambda)$, then
	\[
		Y=\sum_{i=1}^{n}X_{i}\sim\Gamma(\lambda,n)
	\]
	\item The sum of i.i.d. gamma random variables follows a gamma distribution. If $X_{1},\ldots,X_{n}\overset{\ptxt{i.i.d.}}{\sim}\Gamma(\lambda,\alpha)$, then
	\[
		Y=\sum_{i=1}^{n}X_{i}\sim\Gamma(\lambda,n\alpha)
	\]
\end{itemize}

\defn{
	$X$ follows a \u{beta distribution} with parameters $a,b$ ($X\sim\Beta(a,b)$) if, for $0\le{}x\le{}1$,
	\[
		f(x)=\frac{\Gamma(a+b)}{\Gamma(a)\Gamma(b)}x^{a-1}(1-x)^{b-1}
	\]
}

\par\noindent
Observe: When $a=b=1$, $X\sim\Unif[0,1]$.\n

\defn{
	The normal distribution, with $\mu=0$ and $\sigma=1$, is called the \u{standard normal distribution}. $Z\sim{}\Normal(0,1)$ is called a standard normal random variable, with
	\[
		f(z)=\phi(z)=\frac{1}{\sqrt{2\pi}}e^{-\frac{1}{2}z^{2}}\qquad{}F(z)=\Phi(z)=P(Z\le{}z)
	\]
}

\par\noindent
We can standardize any normal random variable: $X\sim\Normal(\mu,\sigma^{2})\Leftrightarrow{}Z=\frac{X-\mu}{\sigma}\sim\Normal(0,1)$. Now, $P(X\le{}x)=\Phi(\frac{x-\mu}{\sigma})$.

\subsection*{Quantiles of the Normal Distribution}

\par\noindent
Notation: $z_{\alpha}$ is the $100(1-\alpha)$ percentile / $1-\alpha$ quantile of the standard normal distribution. Say $X\sim\Normal(\mu,\sigma^{2})$. Let $\eta_{p}$ be the $p$ quantile of $X$. By standardizing, $P(X\le\eta_{p})=\Phi(\frac{\eta_{p}-\mu}{\sigma})$. We then solve $\Phi(\frac{\eta_{p}-\mu}{\sigma})=p$, and get $z_{1-p}=\frac{\eta_{p}-\mu}{\sigma}$, i.e., $\eta_{p}=\mu+\sigma{}z_{1-p}$.\n

\par\noindent
Say $X$ is a random variable with pdf $f_{X}$ and cdf $F_{X}$. Let $Y=g(X)$. We want to find $f_{Y}=f_{g(X)}$.\n

\par\noindent
Approach 1: The direct approach -- work through the cdf. Derive $F_{Y}$ and differentiate.\n

\ex{
	$X\sim\Normal(0,1)$, $Y=X^{2}$. Then $f_{Y}(y)=y^{-1/2}\phi(\sqrt{y})$ for $y\ge{}0$, and $0$ otherwise.\nn
	Proof: For $y\ge{}0$,
	\[
		F_{Y}(y)=P(X^{2}\le{}y)=P(-\sqrt{y}\le{}X\le{}\sqrt{y})=F_{X}(\sqrt{y})-F_{X}(-\sqrt{y})=\Phi(\sqrt{y})-\Phi(-\sqrt{y})
	\]
	\[
		f_{Y}(y)=\frac{d}{dy}F_{Y}(y)=\frac{d}{dy}(\Phi(\sqrt{y})-\Phi(-\sqrt{y}))=\frac{1}{2}\phi(y)y^{-1/2}-(-1\cdot\tfrac{1}{2})\phi(y)y^{-1/2}=\phi(y)y^{-1/2}
	\]
	\proven
}

\ex{
	If $X\sim\Normal(\mu,\sigma^{2})$, $Y=aX+b$, then $Y\sim\Normal(a\mu+b,a^{2}\sigma^{2})$.\n
}

\par\noindent
Approach 2: Chain rule.\n

\thm{
	(Chain Rule) Suppose $f,g$ are differentiable, and $h=g\of{}f$. Then $h'=(g'\of{}f)f'$.\n
}

\par\noindent
Now suppose $g$ is differentiable, and monotonic on an interal $I$, with $f_{X}(x)=0$ outside of $I$. Then
\[
	f_{Y}(y)=f_{X}(g\inv(y))\abs{\frac{dg\inv(y)}{dy}}=f_{X}(g\inv(y))\abs{\frac{1}{g'(g\inv(y))}}
\]
for all $y$ \st{} $\exists{}x\in{}I$ \st{} $y=g(x)$. Otherwise, $f_{Y}(y)=0$.\n

\defn{
	The \u{joint probability mass function} of a pair of discrete random variables $X,Y$ is $p(x,y)=P(X=x\land{}Y=y)$.\n
}

\par\noindent
The joint pmf must satisfy
\begin{itemize}
	\item $p(x,y)\ge{}0$, $\forall(x,y)$
	\item $\sum_{x}\sum_{y}p(x,y)=1$
\end{itemize}

\defn{
	The \u{marginal probability mass function} of $X$, given a joint probability mass function $p$, is $p_{X}(x)=\sum_{y}p(x,y)$.\n
}

\defn{
	$X$ and $Y$ are \u{independent} if $\forall{}(x,y)$, $p(x,y)=p_{X}(x)p_{Y}(y)$.\n
}

\defn{
	The \u{cumulative distribution function} of $X$ and $Y$ is
	\[
		F(x,y)=P(X\le{}x,Y\le{}y)=\sum_{\substack{x_{i}\le{}x\\ y_{i}\le{}y}}p(x_{i},y_{i})
	\]
}

\par\noindent
For $X_{1},\ldots,X_{m}$, the pdf is $p(x_{1},\ldots,x_{m})=P(X_{1}=x_{1},\ldots,X_{m}=x_{m})$.\n
The marginal pdf of $X_{k}$ is
\[
	p_{X_{k}}(x)=\sum_{\mathclap{\substack{x_{1},\ldots,x_{k-1},\\ x_{k+1},\ldots,x_{m}}}}p(x_{1},\ldots,x_{k-1},x,x_{k+1},\ldots,x_{m})
\]
The marginal pdf of $X_{k}$ and $X_{l}$ is (WOLOG $k<l$)
\[
	p_{X_{k},X_{l}}(s,t)=P(X_{k}=s,X_{l}=t)=\sum_{\mathclap{\substack{x_{1},\ldots{}x_{k-1},\\ x_{k+1},\ldots,x_{l-1},\\ x_{l+1},\ldots,x_{m}}}}p(x_{1},\ldots,x_{k-1},s,x_{k+1},\ldots,x_{l-1},t,x_{l+1},\ldots,x_{m})
\]

\ex{
	A fair coin is tossed $3$ times. $X$ is $1$ if the first toss is heads, and $0$ if it's tails. $Y$ is the total number of heads. Then the sample space is
	\[
		\set{(H,H,H),(H,H,T),(H,T,H),(H,T,T),(T,H,H),(T,H,T),(T,T,H),(T,T,T)}
	\]
	The pdf of $(X,Y)$ is
	\[
		\begin{array}{cc|cccc}
			& & \multicolumn{4}{c}{Y}\\
			& p(x,y) & 0 & 1 & 2 & 3 \\ \hline
			& & \\[-10pt]
			\multirow{2}{*}{X} & 0 & \frac{1}{8} & \frac{2}{8} & \frac{1}{8} & 0\\
			& & \\[-10pt]
			& 1 & 0 & \frac{1}{8} & \frac{2}{8} & \frac{1}{8}
		\end{array}
	\]
	$P(X=0,Y=2)=\frac{1}{8}$\n
	$P(Y=2)=P(X=0,Y=2)+P(X=1,Y=2)=\frac{3}{8}$\n
	$F_{X,Y}(0,2)=P(X\le{}0,Y\le{}2)=\frac{1}{8}+\frac{2}{8}+\frac{1}{8}=\frac{1}{2}$\n
	The marginal cdf of $Y$ is
	\[
		\begin{array}{c|cccc}
			Y & 0 & 1 & 2 & 3\\ \hline
			& & \\[-10pt]
			\ptxt{pdf} & \frac{1}{8} & \frac{3}{8} & \frac{3}{8} & \frac{1}{8}\\
			& & \\[-10pt]
			\ptxt{cdf} & \frac{1}{8} & \frac{1}{2} & \frac{7}{8} & 1
		\end{array}
	\]
}

\end{document}